\uuid{cjx3}
\titre{Tangente}
\auteur{bodin ; hanani ; mzari}
\organisation{exo7}
\licence{CC BY-NC-SA 4.0}
\datecreate{2020-01-20}
\niveau{L1}
\module{Analyse}
\chapitre{Dérivabilité des fonctions réelles}
\sousChapitre{Autre}
\difficulte{1}
\theme{Dérivées, tangente}
\exo7qcmid{lille-1.derivees.002}

\contenu{
\begin{question}
\qclassification{124.00}
\qtags{motcle=tangente}
Etant donné que $\displaystyle f(3)=1$ et $f'(3)=5$. Une équation de la tangente à $\mathcal{C}_f$ au point $(3,1)$ est :
\begin{answers}
    \bad{$y=1(x-3)+5=x+2$}
    \bad{$y=1(x-3)-5=x-8$}
    \bad{$y=5(x-3)-1=5x-16$}
    \good{$y=5(x-3)+1=5x-14$}
\end{answers}
\begin{explanations}
On applique la formule du cours $\displaystyle y=f'(3)(x-3)+f(3)=5(x-3)+1$.
\end{explanations}
\end{question}
}
